% ============================================================
% Section 3: Advantages & Limitations
% ============================================================
\section{Advantages \& Limitations}

\begin{frame}{Advantages}
    \begin{enumerate}
        \item \textbf{No predefined cluster count}
        \begin{itemize}
            \item Automatically discovers number of clusters
            \item Unlike K-means which requires $k$
        \end{itemize}

        \vspace{0.3cm}

        \item \textbf{Flexible cluster shapes}
        \begin{itemize}
            \item Can find non-spherical clusters
            \item Follows natural density contours
        \end{itemize}

        \vspace{0.3cm}

        \item \textbf{Single parameter: bandwidth $h$}
        \begin{itemize}
            \item Controls granularity of clustering
            \item Easier to tune than multiple parameters
        \end{itemize}

        \vspace{0.3cm}

        \item \textbf{Robust to initialization}
        \begin{itemize}
            \item Deterministic convergence to modes
            \item No random initialization sensitivity
        \end{itemize}
    \end{enumerate}
\end{frame}

\begin{frame}{Limitations (and Why They Matter)}
    \begin{enumerate}
        \item \textbf{Computational complexity: $O(n^2 \cdot T)$}
        \begin{itemize}
            \item[\textcolor{customred}{Why?}] Each point computes distance to all others at each iteration
            \item[\textcolor{customred}{Impact:}] Slow for large datasets ($n > 10{,}000$)
            \item[\textcolor{orange}{Mitigation:}] Use spatial indexing (KD-trees), binning, or sampling
        \end{itemize}

        \vspace{0.3cm}

        \item \textbf{Bandwidth selection is critical}
        \begin{itemize}
            \item[\textcolor{customred}{Why?}] Too small $h$ → many tiny clusters; too large $h$ → all points in one cluster
            \item[\textcolor{customred}{Impact:}] Results highly sensitive to $h$
            \item[\textcolor{orange}{Mitigation:}] Cross-validation, adaptive bandwidth, or domain knowledge
        \end{itemize}

        \vspace{0.3cm}

        \item \textbf{Curse of dimensionality}
        \begin{itemize}
            \item[\textcolor{customred}{Why?}] Density estimation becomes unreliable in high dimensions
            \item[\textcolor{customred}{Impact:}] Poor performance for $d > 10$
            \item[\textcolor{orange}{Mitigation:}] Dimensionality reduction (PCA, t-SNE) first
        \end{itemize}
    \end{enumerate}
\end{frame}

\begin{frame}{When to Use Mean Shift?}
    \begin{columns}[T]
        \begin{column}{0.5\textwidth}
            \textbf{\textcolor{green!60!black}{Good Scenarios:}}
            \begin{itemize}
                \item Unknown number of clusters
                \item Non-spherical cluster shapes
                \item Small to medium datasets
                \item Image segmentation
                \item Object tracking in video
                \item Mode-seeking tasks
            \end{itemize}
        \end{column}

        \begin{column}{0.5\textwidth}
            \textbf{\textcolor{red}{Avoid When:}}
            \begin{itemize}
                \item Very large datasets ($n > 50{,}000$)
                \item High-dimensional data ($d > 10$)
                \item Real-time processing needed
                \item Clusters have similar densities
                \item Well-separated spherical clusters (use K-means instead)
            \end{itemize}
        \end{column}
    \end{columns}
\end{frame}
